\documentclass[10pt,a4paper]{article}
\usepackage{amsmath,amssymb,amsthm,graphicx}
\usepackage{amsmath,amsfonts,amssymb,amsthm,enumerate,hyperref,ytableau,authblk,listings}
\usepackage[capitalize]{cleveref}
\newtheorem{theorem}{Theorem}[section]
\newtheorem{lemma}[theorem]{Lemma}
\newtheorem{corollary}[theorem]{Corollary}
\newtheorem{conjecture}[theorem]{Conjecture}

\theoremstyle{definition}

\newtheorem{definition}[theorem]{Definition}
\providecommand{\coset}[2]{[#1 \mapsto #2]}
\providecommand{\dcoset}[2]{[#1 \leftrightarrow #2]}
\providecommand{\id}{\mathsf{id}}
\providecommand{\cM}{\mathcal{M}}
\title{An m-covered family gives every member a certificate of size m plus one}
\author{Gilad Chase, Neta Dafni, Yuval Filmus, Nathan Lindzey}
\date{arxiv:2210.00245v2; Creative Commons Attribution 4.0 International}
\begin{document}
\maketitle
\noindent\textit{Editorial note.} The counted statement is the author's labelled environment at source0.tex lines 581-585 and its proof environment opens at line 586; the pairing is by adjacency and the retained proof range is the whole environment.\par\medskip
\noindent\textit{Editorial note.} Bibliography: the article's own compiled bibliography file (main.bbl, a complete thebibliography with 29 entries) is bundled unchanged; no entry is edited or re-rendered.\par\medskip
\noindent\textit{Editorial note.} Retained-licence form: four-key reuse record in a task-directory evidence copy; no status key and no acceptance record are written.\par\medskip
\noindent\textit{Editorial note.} Metadata provenance: source\_version from the source's licence evidence; authors from the authors\_index key of its source record, title from the title key.\par\medskip
\noindent\textit{Editorial note.} No label-number dump exists for this source and no counter values are imposed, so identity is the author label.\par\medskip
\noindent\textit{Editorial scope.} One original result of the pinned author manuscript, counted once: the certificate bound for an m-covered family (author label lem:structure-sym).\par\medskip
\noindent\textit{Allowed standard prerequisites (editorial).}\begin{itemize}
\item m-covered family. The family is contained in the union of m cosets, each prescribing one image of a permutation.. It is the hypothesis.
\item Certificate of a member. A set of images prescribed for that member; the size bound is the conclusion.. It is what the statement bounds.
\item Degree reduction. Reducing a function by fixing the coordinates of a coset gives a function of lower degree on the remaining coordinates.. It is the tool that lets the degree-one classification be applied.
\end{itemize}\medskip
\begin{lemma} \label{lem:degree1-cert-sym}
If $f\colon S_n \to \{0,1\}$ has degree at most~$1$ then $C(f) \leq 1$.
\end{lemma}

\begin{lemma} \label{lem:degree-reduction-sym}
Suppose that $f\colon S_n \to \{0,1\}$ is the characteristic function of a family which is a subset of the coset $\coset{i}{j}$. Then $\deg f|_{\coset{i}{j}} \leq \max( \deg f - 1, 0 )$.
\end{lemma}

\begin{lemma} \label{lem:degree1-cert-pms}
If $f\colon \cM_{2n} \to \{0,1\}$ has degree at most~$1$ then $C(f) \leq 2$.
\end{lemma}

\begin{lemma} \label{lem:structure-sym}
Let $\mathcal{F} \subseteq S_n$ be an $m$-covered family, say $\mathcal{F} \subseteq \coset{i_1}{j_1} \cup \cdots \cup \coset{i_m}{j_m}$. Suppose that its characteristic function $f$ has degree at most~$2$.

Every $\alpha$ has a certificate of the form $\{(i_1,\alpha(i_1)),\ldots,(i_m,\alpha(i_m)),(i,j)\}$, and in particular, $C(f) \leq m + 1$.
\end{lemma}

\begin{proof}
We will prove by induction on $m$ that if $\mathcal{F} \subseteq S_n$ is $m$-covered, say $\mathcal{F} \subseteq \coset{i_1}{j_1} \cup \cdots \cup \coset{i_m}{j_m}$, and its characteristic function $f$ has degree at most~$2$, then for every distinct $k_1,\ldots,k_m \in [n]$, the restriction of $f$ to $\coset{i_1}{k_1} \cap \cdots \cap \coset{i_m}{k_m}$ has degree at most~$1$. By \Cref{lem:degree1-cert-sym}, this restriction has certificate complexity at most~$1$. For any $\alpha \in S_n$, choosing $k_1 = \alpha(i_1),\ldots,k_m = \alpha(i_m)$, we deduce that $\alpha$ has a certificate of the required form.

We start with the base case $m = 1$. Let $\mathcal{F} \subseteq S_n$ be such that $\mathcal{F} \subseteq \coset{i_1}{j_1}$, and suppose that its characteristic function $f$ has degree at most~$2$. \Cref{lem:degree-reduction-sym} directly implies that the restriction of $f$ to $\coset{i_1}{j_1}$ has degree at most~$1$, as required.

Suppose now that $m \geq 2$. Let $\mathcal{F} \subseteq S_n$ be such that $\mathcal{F} \subseteq \coset{i_1}{j_1} \cup \cdots \cup \coset{i_m}{j_m}$, and suppose that its characteristic function $f$ has degree at most~$2$. Consider any distinct $k_1,\ldots,k_m \in [n]$. We need to show that the restriction of $f$ to $\coset{i_1}{k_1} \cap \cdots \cap \coset{i_m}{k_m}$ has degree at most~$1$.

Suppose first that $k_s \neq j_s$ for some $s \in [m]$, say $k_m \neq j_m$. The restriction of $\mathcal{F}$ to $\coset{i_m}{k_m}$ is contained in $\coset{i_1}{j_1} \cup \cdots \cup \coset{i_{m-1}}{j_{m-1}}$, and its characteristic function has degree at most~$2$. Therefore the induction hypothesis implies that the restriction of $f$ to $\coset{i_1}{k_1} \cap \cdots \cap \coset{i_m}{k_m}$ has degree at most~$1$.

It remains to show that the restriction of $f$ to $\coset{i_1}{j_1} \cap \cdots \cap \coset{i_m}{j_m}$ has degree at most~$1$. To see this, consider the function
\[
 g := x_{i_1j_1} \cdots x_{i_mj_m} f = f - \sum_{\substack{k_1,\ldots,k_m \\ (k_1,\ldots,k_m) \neq (j_1,\ldots,j_m)}} x_{i_1k_1} \cdots x_{i_mk_m} f.
\]
We showed above that the restriction of $f$ to $\coset{i_1}{k_1} \cap \cdots \cap \coset{i_m}{k_m}$ has degree at most~$1$ whenever $(k_1,\ldots,k_m) \neq (j_1,\ldots,j_m)$. Therefore every term in the sum on the right-hand side has degree at most $m+1$, since it can be represented as $x_{i_1k_1} \cdots x_{i_mk_m}$ times the linear polynomial representing the restriction of $f$ to $\coset{i_1}{k_1} \cap \cdots \cap \coset{i_m}{k_m}$. Consequently, $\deg g \leq \max(\deg f, m+1) = m+1$.

By construction, $g$ is the characteristic function of a family which is contained in the coset intersection $\mathcal{I} = \coset{i_1}{j_1} \cap \cdots \cap \coset{i_m}{j_m}$. Applying \Cref{lem:degree-reduction-sym} repeatedly, we conclude that the restriction of $g$ to $\mathcal{I}$ has degree at most $\max(m+1-m,0) = 1$. Since the restriction of $f$ to $\mathcal{I}$ is the same as the restriction of $g$ to $\mathcal{I}$, this completes the proof of the induction step.
\end{proof}

\newcommand{\etalchar}[1]{$^{#1}$}
\begin{thebibliography}{BMRR21}

\bibitem[AK97]{AK1}
Rudolf Ahlswede and Levon~H. Khachatrian.
\newblock The complete intersection theorem for systems of finite sets.
\newblock {\em European J. Combin.}, 18(2):125--136, 1997.

\bibitem[AK99]{AK2}
Rudolf Ahlswede and Levon~H. Khachatrian.
\newblock A pushing-pulling method: new proofs of intersection theorems.
\newblock {\em Combinatorica}, 19(1):1--15, 1999.

\bibitem[BdW02]{BuhrmandeWolf}
Harry Buhrman and Ronald de~Wolf.
\newblock Complexity measures and decision tree complexity: a survey.
\newblock {\em Theoret. Comput. Sci.}, 288(1):21--43, 2002.
\newblock Complexity and logic (Vienna, 1998).

\bibitem[BMRR21]{BMRR}
Angelot Behajaina, Roghayeh Maleki, Aina~Toky Rasoamanana, and A.~Sarobidy
  Razafimahatratra.
\newblock 3-setwise intersecting families of the symmetric group.
\newblock {\em Discrete Math.}, 344(8):Paper No. 112467, 15, 2021.

\bibitem[CK03]{CameronKu}
Peter~J. Cameron and C.~Y. Ku.
\newblock Intersecting families of permutations.
\newblock {\em European J. Combin.}, 24(7):881--890, 2003.

\bibitem[CSST08]{CSST08}
Tullio Ceccherini-Silberstein, Fabio Scarabotti, and Filippo Tolli.
\newblock {\em Harmonic analysis on finite groups}, volume 108 of {\em
  Cambridge Studies in Advanced Mathematics}.
\newblock Cambridge University Press, Cambridge, 2008.
\newblock Representation theory, Gelfand pairs and Markov chains.

\bibitem[Daf22]{DafniThesis}
Neta Dafni.
\newblock Complexity measures on the symmetric group and beyond.
\newblock Master's thesis, Technion --- Israel Institute of Technology, 2022.

\bibitem[DFL{\etalchar{+}}21]{DFLLV}
Neta Dafni, Yuval Filmus, Noam Lifshitz, Nathan Lindzey, and Marc Vinyals.
\newblock Complexity measures on the symmetric group and beyond.
\newblock In {\em 12th {I}nnovations in {T}heoretical {C}omputer {S}cience
  {C}onference}, volume 185 of {\em LIPIcs. Leibniz Int. Proc. Inform.}, pages
  Art. No. 87, 5. Schloss Dagstuhl. Leibniz-Zent. Inform., Wadern, 2021.

\bibitem[EFP11]{EFP}
David Ellis, Ehud Friedgut, and Haran Pilpel.
\newblock Intersecting families of permutations.
\newblock {\em J. Amer. Math. Soc.}, 24(3):649--682, 2011.

\bibitem[EKR61]{EKR}
P.~Erd\H{o}s, Chao Ko, and R.~Rado.
\newblock Intersection theorems for systems of finite sets.
\newblock {\em Quart. J. Math. Oxford Ser. (2)}, 12:313--320, 1961.

\bibitem[Ell12a]{EllisCameronKu}
David Ellis.
\newblock A proof of the {C}ameron-{K}u conjecture.
\newblock {\em J. Lond. Math. Soc. (2)}, 85(1):165--190, 2012.

\bibitem[Ell12b]{EllisSetwise}
David Ellis.
\newblock Setwise intersecting families of permutations.
\newblock {\em J. Combin. Theory Ser. A}, 119(4):825--849, 2012.

\bibitem[FD77]{DezaFrankl}
P\'{e}ter Frankl and Mikhail Deza.
\newblock On the maximum number of permutations with given maximal or minimal
  distance.
\newblock {\em J. Combinatorial Theory Ser. A}, 22(3):352--360, 1977.

\bibitem[FHKL16]{FHKL}
Yuval Filmus, Hamed Hatami, Nathan Keller, and Noam Lifshitz.
\newblock On the sum of the {$L_1$} influences of bounded functions.
\newblock {\em Israel J. Math.}, 214(1):167--192, 2016.

\bibitem[Fil17]{FilmusComment}
Yuval Filmus.
\newblock A comment on intersecting families of permutations.
\newblock arXiv:1706.10146, 2017.

\bibitem[FMS21]{FMS}
Shaun Fallat, Karen Meagher, and Mahsa~N. Shirazi.
\newblock The {E}rd{\H{o}}s-{K}o-{R}ado theorem for 2-intersecting families of
  perfect matchings.
\newblock {\em Algebr. Comb.}, 4(4):575--598, 2021.

\bibitem[GM09]{GodsilMeagher}
Chris Godsil and Karen Meagher.
\newblock A new proof of the {E}rd{\H{o}}s-{K}o-{R}ado theorem for intersecting
  families of permutations.
\newblock {\em European J. Combin.}, 30(2):404--414, 2009.

\bibitem[GM16]{GodsilMeagherBook}
Chris Godsil and Karen Meagher.
\newblock {\em Erd\H{o}s-{K}o-{R}ado theorems: algebraic approaches}, volume
  149 of {\em Cambridge Studies in Advanced Mathematics}.
\newblock Cambridge University Press, Cambridge, 2016.

\bibitem[GM17]{GodsilMeagherPM}
Chris Godsil and Karen Meagher.
\newblock An algebraic proof of the {E}rd{\H{o}}s-{K}o-{R}ado theorem for
  intersecting families of perfect matchings.
\newblock {\em Ars Math. Contemp.}, 12(2):205--217, 2017.

\bibitem[Lin17]{Lindzey17}
Nathan Lindzey.
\newblock Erd{\H{o}}s-{K}o-{R}ado for perfect matchings.
\newblock {\em European J. Combin.}, 65:130--142, 2017.

\bibitem[Lin18a]{Lindzey18}
Nathan Lindzey.
\newblock Intersecting families of permutations of perfect matchints.
\newblock arXiv:1811.06160, 2018.

\bibitem[Lin18b]{LindzeyThesis}
Nathan Lindzey.
\newblock {\em Matchings and Representation Theory}.
\newblock PhD thesis, University of Waterloo, 2018.

\bibitem[LM04]{LaroseMalvenuto}
Benoit Larose and Claudia Malvenuto.
\newblock Stable sets of maximal size in {K}neser-type graphs.
\newblock {\em European J. Combin.}, 25(5):657--673, 2004.

\bibitem[MM05]{MeagherMoura}
Karen Meagher and Lucia Moura.
\newblock Erd{\H{o}}s-{K}o-{R}ado theorems for uniform set-partition systems.
\newblock {\em Electron. J. Combin.}, 12:Research Paper 40, 12, 2005.

\bibitem[MR21]{MR21}
Karen Meagher and A.~Sarobidy Razafimahatratra.
\newblock The {E}rd{\H{o}}s-{K}o-{R}ado theorem for 2-pointwise and 2-setwise
  intersecting permutations.
\newblock {\em Electron. J. Combin.}, 28(4):Paper No. 4.10, 21, 2021.

\bibitem[N{\"O}03]{Cliquer}
Sampo Niskanen and Patric R.~J. {\"O}sterg{\aa}rd.
\newblock Cliquer user's guide, version 1.0.
\newblock Technical Report T48, Communications Laboratory, Helsinki University
  of Technology, Espoo, Finland, 2003.

\bibitem[O'D14]{ODonnell}
Ryan O'Donnell.
\newblock {\em Analysis of {B}oolean functions}.
\newblock Cambridge University Press, New York, 2014.

\bibitem[Sag01]{Sagan}
Bruce~E. Sagan.
\newblock {\em The symmetric group}, volume 203 of {\em Graduate Texts in
  Mathematics}.
\newblock Springer-Verlag, New York, second edition, 2001.
\newblock Representations, combinatorial algorithms, and symmetric functions.

\bibitem[{Sag}22]{sagemath}
{Sage Developers}.
\newblock {\em {S}ageMath, the {S}age {M}athematics {S}oftware {S}ystem
  ({V}ersion 9.5)}, 2022.
\newblock {\tt https://www.sagemath.org}.

\end{thebibliography}
\end{document}
